\documentclass[10pt,aspectratio=169]{beamer}
\input{../../_en_preambulo_beamer}% Comandos y macros estadísticas compartidas para las presentaciones Beamer en inglés (Metropolis)

% Conjuntos numéricos básicos
\newcommand{\R}{\mathbb{R}}
\newcommand{\N}{\mathbb{N}}
\newcommand{\Q}{\mathbb{Q}}
\newcommand{\Z}{\mathbb{Z}}
\newcommand{\set}[1]{\left\{ #1 \right\}}
\newcommand{\abs}[1]{\left|#1\right|}
\newcommand{\norm}[1]{\left\| #1 \right\|}

% Comandos estadísticos y probabilísticos del curso
\newcommand{\Var}{\operatorname{Var}}
\newcommand{\cov}{\operatorname{Cov}}
\newcommand{\corr}{\rho}
\newcommand{\comb}[2]{\begin{pmatrix} #1 \\ #2 \end{pmatrix}}
\newcommand{\s}{\sigma}
\newcommand{\particion}{\mathfrak{P}}
\newcommand{\card}[1]{\#\left( #1 \right)}
\newcommand{\seq}[3]{\set{#1}_{#2}^{#3}}
\newcommand{\E}{\mathbb{E}}
\newcommand{\Prob}{\mathbb{P}}

% Entornos pedagógicos y cajas en inglés académico natural (soportando tanto nombres en inglés como en español)
\newenvironment{discussion}{%
  \begin{alertblock}{Discussion Question}%
}{%
  \end{alertblock}%
}
\newenvironment{discusion}{%
  \begin{alertblock}{Discussion Question}%
}{%
  \end{alertblock}%
}

\newenvironment{reflection}{%
  \begin{alertblock}{Classroom Reflection}%
}{%
  \end{alertblock}%
}
\newenvironment{reflexion}{%
  \begin{alertblock}{Classroom Reflection}%
}{%
  \end{alertblock}%
}

\newenvironment{paradox}{%
  \begin{alertblock}{Exploratory Paradox}%
}{%
  \end{alertblock}%
}
\newenvironment{paradoja}{%
  \begin{alertblock}{Exploratory Paradox}%
}{%
  \end{alertblock}%
}

\newenvironment{motivation}{%
  \begin{exampleblock}{Motivation: Data Science in Practice}%
}{%
  \end{exampleblock}%
}
\newenvironment{motivacion}{%
  \begin{exampleblock}{Motivation: Data Science in Practice}%
}{%
  \end{exampleblock}%
}

\newenvironment{quickchallenge}{%
  \begin{block}{Quick Team Challenge}%
}{%
  \end{block}%
}
\newenvironment{retoclase}{%
  \begin{block}{Quick Team Challenge}%
}{%
  \end{block}%
}

\title[Experimental strategies]{Section 08.01: Experimental Strategies}\subtitle{Design of Experiments and ANOVA}
\author[J. Castillo Colmenares]{Juliho Castillo Colmenares}\institute{Tecnológico de Monterrey}\date{}
\begin{document}
\begin{frame}[plain]\titlepage\end{frame}
\begin{frame}{Roadmap}\small\begin{enumerate}\item Motivation and terminology.\item Fisher's principles.\item Planning and analysis.\end{enumerate}\end{frame}
\begin{frame}{Source and scope}\small This deck follows the experimental-strategies section and bridges Design of Experiments (DoE) with ANOVA.\begin{block}{Guiding question}How can a design distinguish causal effects from noise?\end{block}\end{frame}
\begin{frame}{Why experiment}\small An observational study can mix treatments with confounders. Controlled assignment compares scenarios under a reproducible protocol.\end{frame}
\begin{frame}{Inflation from pairwise tests}\small Comparing $k$ groups pairwise creates many false-positive opportunities. For $m$ independent comparisons,\[\alpha_{\mathrm{global}}=1-(1-\alpha)^m.\]\end{frame}
\begin{frame}{Essential terminology}\small\begin{itemize}\item Experimental unit: entity receiving a treatment.\item Factor: manipulated variable.\item Treatment: concrete combination of levels.\item Response: measured variable.\end{itemize}\end{frame}
\begin{frame}{Replication}\small Repeating a treatment on independent units estimates error variance and reduces the standard error of a treatment mean. Without replication, effect and noise cannot be separated.\end{frame}
\begin{frame}{Randomization}\small Randomly assigning treatments prevents systematic unit differences from being confused with factor effects and supports the probability model.\end{frame}
\begin{frame}{Blocking}\small Grouping relatively homogeneous units before assignment controls a known variability source and leaves a smaller residual error.\end{frame}
\begin{frame}{Experimental protocol}\small\begin{enumerate}\item Define question and response.\item Choose factors and levels.\item Set units, replication, and blocks.\item Randomize assignment.\end{enumerate}\end{frame}
\begin{frame}{Bridge to ANOVA}\small A well-planned design produces comparable means and an internal error estimate. ANOVA tests several treatments in one model.\end{frame}
\begin{frame}{Design matrix}\small Record each unit, factor levels, block when present, randomized order, and observed response. This provides traceability.\end{frame}
\begin{frame}{Application: A/B/n test}\small Compare several algorithm configurations. Each configuration is a treatment; execution time is the response.\end{frame}
\begin{frame}{Application: family-wise risk}\small With $k=5$ treatments there are $m=10$ pairs. If every comparison uses $\alpha=0.05$,\[1-(0.95)^{10}\approx0.4013.\]The global risk exceeds the nominal level.\end{frame}
\begin{frame}{Application: replication}\small Run each configuration on several independent units to estimate time variability and distinguish stable differences from fluctuation.\end{frame}
\begin{frame}{Application: randomization}\small Randomize execution order so warm-up, time of day, or server load cannot be confounded with treatment.\end{frame}
\begin{frame}{Application: blocking}\small If servers belong to different generations, use generation as a block and assign configurations within each block.\end{frame}
\begin{frame}{Application: analysis}\small The final plan combines treatment, replication, randomization, and blocking; then apply ANOVA and adjusted comparisons when appropriate.\end{frame}
\begin{frame}{Interpretation}\small Causal inference depends on design as well as calculation. A $p$-value cannot repair biased assignment or a poorly defined response.\end{frame}
\begin{frame}{Communication}\small Report units, factors, levels, replicates, blocks, randomization rule, and response so the analysis can be reproduced.\end{frame}
\begin{frame}{Synthesis}\small\begin{itemize}\item Replication estimates error.\item Randomization protects against confounding.\item Blocking controls known heterogeneity.\item ANOVA compares treatments.\end{itemize}\end{frame}
\begin{frame}{Chapter connection}\small The next sections develop one-factor ANOVA, fixed effects, blocked designs, and factorial designs.\end{frame}
\end{document}
